\chapter{星系形态分类案例}

\section{案例目标}

这一章把 Part III 里的分类、特征工程、模型评估和可信度判断，放回一个最经典也最容易被误解的天文图像任务：星系形态分类。目标是让读者看到，一个看似“很适合深度学习”的问题，真正做起来其实包含了大量前处理和标签解释工作。

\begin{enumerate}
    \item 理解星系形态分类究竟在回答什么科学问题；
    \item 学会从图像 cutout、尺度统一和简单质量控制开始组织数据；
    \item 了解传统手工特征与卷积网络各自适合什么阶段；
    \item 认识类别不平衡、投影效应和标签噪声对分类结果的影响；
    \item 把模型输出重新放回形态学和观测条件中解释。
\end{enumerate}

\section{问题背景：形态分类为什么重要}

星系形态分类并不只是“给图片贴标签”。在天体物理上，它经常和更深的科学问题连在一起，例如：

\begin{itemize}
    \item 星系形成和并合历史；
    \item 恒星形成活跃度与结构之间的关系；
    \item 环境对盘星系、棒旋星系和椭圆星系比例的影响；
    \item 大样本巡天里形态统计和演化趋势的测量。
\end{itemize}

也就是说，形态标签本身往往只是一个中间量。真正的目标，是借助这些标签去问更大的科学问题。

\section{数据：从图像 cutout 到标签}

这个案例面向的是 Galaxy Zoo 风格或巡天缩略图风格的小型教学数据集。最典型的数据组织形式是：

\begin{itemize}
    \item 每个对象有一张对齐后的图像 cutout；
    \item 同时配有一个教学用形态标签，例如 \texttt{spiral}、\texttt{elliptical}、\texttt{edge\_on}；
    \item 附带少量观测信息，例如信噪比、视宁度或红移区间。
\end{itemize}

真正项目里，数据看上去往往没有想象中整洁。常见问题包括：

\begin{itemize}
    \item 图像尺度不一致；
    \item 星系中心不在 cutout 正中；
    \item 低信噪比导致旋臂几乎不可见；
    \item 邻近亮源或背景结构污染轮廓；
    \item 人工标签本身存在分歧。
\end{itemize}

所以，形态分类的第一课同样不是选模型，而是先想清楚图像到底可不可比。

\section{第一步：先把图像组织成可比较状态}

对教学案例来说，最基本的前处理通常包括：

\begin{itemize}
    \item 统一 cutout 尺寸；
    \item 保证目标大致居中；
    \item 规范化像素强度范围；
    \item 记录明显的低质量样本。
\end{itemize}

\begin{examplebox}[一个最小图像预处理思路]
\begin{lstlisting}[style=pythonstyle]
image = crop_center(raw_image)
image = rescale_to_fixed_size(image)
image = normalize_pixels(image)
\end{lstlisting}
\end{examplebox}

这一部分看起来“不是 AI”，但它往往直接决定后面的分类结果。一个模型如果主要学到的是背景亮度、噪声模式或 cutout 边缘，而不是星系结构，那么它的高准确率也没有太大意义。

\section{第二步：传统特征为什么仍然值得保留}

今天一提到图像分类，很多人第一反应都是 CNN。但在教学和真实科研早期阶段，传统手工特征仍然非常有价值。例如：

\begin{itemize}
    \item 轴比、集中度、非对称度；
    \item 表面亮度分布；
    \item 纹理统计；
    \item 简单的径向轮廓特征。
\end{itemize}

这些特征的好处是：

\begin{itemize}
    \item 解释性强；
    \item 样本较小时更稳；
    \item 有助于判断“图像里到底是哪类结构在起作用”。
\end{itemize}

因此，这一章非常适合把“传统特征 baseline”和“卷积网络 baseline”并列来看，而不是直接跳过前者。

\section{第三步：什么时候 CNN 真正有优势}

当数据预处理比较稳定、样本量也开始够用时，卷积神经网络会自然展现优势。原因主要有三点：

\begin{itemize}
    \item 它能直接从局部结构中学习旋臂、核、棒和不规则轮廓；
    \item 它不必依赖手工定义每个形态指标；
    \item 迁移学习可以显著降低从头训练的门槛。
\end{itemize}

但要记住：CNN 的优势建立在数据条件也比较靠谱的前提下。如果标签本身噪声很大、低信噪比样本太多、类间边界本来就模糊，那么复杂模型不一定能解决核心问题。

\section{第四步：投影效应和类别边界是这个案例的难点}

星系形态分类最棘手的地方之一，是有些“错分”未必真的是模型错了。例如：

\begin{itemize}
    \item 边缘朝向的盘星系可能看起来像细长椭圆体；
    \item 低信噪比螺旋星系可能只剩一个模糊亮斑；
    \item 并合或扰动系统本来就不适合硬塞进简单类别；
    \item 不同标注者对同一张图的判断可能并不一致。
\end{itemize}

这意味着你在看混淆矩阵时，要特别小心：某些混淆反映的不是模型不行，而是类别本身就带有连续过渡和人为定义。

\section{第五步：类别不平衡和标签噪声不能忽视}

在很多真实样本里，最常见的类型通常数量最多，而真正罕见但科学上有趣的类型数量很少。于是你常会同时面对：

\begin{itemize}
    \item 类别不平衡；
    \item 少数类 recall 很差；
    \item 某些类别训练样本过少；
    \item 标签之间的投票分歧。
\end{itemize}

这时仅仅报告 overall accuracy 基本是不够的。你至少还应关心：

\begin{itemize}
    \item 每类召回率；
    \item 哪些类别最容易互相混淆；
    \item 低质量样本是否集中在某些错分区域；
    \item 是否需要把“高不确定度样本”单独交回人工复核。
\end{itemize}

\section{为什么这个案例特别适合比较传统方法和现代 AI}

星系形态分类有一个很好的教学特点：它天然允许多条方法线并存。

\begin{itemize}
    \item 传统图像特征 + 经典分类器；
    \item CNN 从像素直接学习；
    \item 迁移学习作为低成本深度学习方案；
    \item 人工复核作为高不确定度样本的闭环。
\end{itemize}

因此，这一章很适合训练学生的工程判断：不是默认“新方法一定更好”，而是结合样本规模、标签质量、算力和科学目标选择恰当复杂度。

\begin{warnbox}[星系形态分类中的常见误区]
\begin{itemize}
    \item 没有处理尺度和居中问题，就直接喂图像给模型；
    \item 用总体准确率掩盖少数类几乎学不到的事实；
    \item 把投影效应导致的边界模糊误当作单纯模型错误；
    \item 忽视人工标签本身的不确定性；
    \item 在样本很小的时候盲目追求从头训练深网络。
\end{itemize}
\end{warnbox}

\section{AI 助手如何帮助你}

在这个案例里，AI 助手很适合帮助你：

\begin{itemize}
    \item 梳理图像前处理和数据增强的实验方案；
    \item 帮你整理混淆矩阵和错分图像复查记录；
    \item 把“传统特征 baseline”和“CNN baseline”写成结构清楚的对照实验说明。
\end{itemize}

但到底是标签有问题、图像质量有问题，还是模型真的没学到结构；以及某个高不确定度对象是否值得单独研究，仍然需要你自己判断。

\section{本章小结}

星系形态分类案例提醒我们：图像 AI 任务真正困难的部分，常常不是网络结构本身，而是图像尺度、标签质量、投影效应和科学定义边界。只有把这些问题一起放进工作流里，分类结果才会真正变得可信。
